% !TEX root = ../main.tex
%-----------------------------------------------------------------------
%  Title block, the question, and the experiment design.
%  \input by ../main.tex -- not compiled on its own.
%-----------------------------------------------------------------------
\begin{center}
{\large\bfseries A Vertical-Mixing Perturbation Ensemble\\[2pt]
 for Ocean Adjoint Sensitivities}\\[3pt]
{\small Tanvir Shahriar \quad$\cdot$\quad DINO configuration, MITgcm
 \quad$\cdot$\quad \today}
\end{center}

\hd{The question.}
An adjoint run gives the sensitivity of the basin's overturning transport to
everything the model depends on --- initial state, surface forcing, mixing
parameters --- in a single backward integration. What it does not tell us is
whether that answer survives a change in the model itself. Every sensitivity map
we have was computed at one setting of the vertical mixing, and whether the
patterns hold at other settings is not known for this configuration.

The experiment asks it directly. Seven adjoint runs, identical to a completed
reference run in every respect except the vertical diffusivity \kv, which is
varied over a factor of $128$.

\hd{Either answer is worth having.}
This is the property that makes the experiment worth running before anything is
built on top of it.

\smallskip
\begin{tabularx}{\textwidth}{@{}Y{0.42}Y{1.58}@{}}
\toprule
\textbf{If the patterns barely move} & Published sensitivity maps are robust to
  the mixing uncertainty that every ocean model carries. That is a result in its
  own right, and it has nothing to do with machine learning\\
\textbf{If they move substantially} & Mixing is a first-order control on
  sensitivity, and anything that tries to generalise or emulate these maps has to
  treat it as an input rather than a constant\\
\bottomrule
\end{tabularx}

\hd{The design.}
Each member starts from the same spin-up state, has its diffusivity changed and
nothing else, is given ten years to re-equilibrate, and is then run through the
same five-year adjoint window as the reference run:

\smallskip
\begin{center}
\footnotesize
\begin{verbatim}
member    :  2170 --[10 yr, perturbed k]--> 2180 --[5 yr adjoint]--> 2185
reference :  2170 --[10 yr, reference k]--> 2180 --[5 yr adjoint]--> 2185
\end{verbatim}
\end{center}
\smallskip

\hd{Why the perturbation leg sits before the reference window.}
This is the one design choice worth pausing on, and it is what makes the
experiment cheap. The spin-up already ran through that decade continuously, so
its own trajectory is the unperturbed counterpart of every member's leg, and the
completed reference run becomes a valid control at no additional cost. Both
branches begin from an identical state, run for the same duration, and are
evaluated over the same window --- which leaves diffusivity as the only
difference between them.

Re-equilibrating forward instead would place the member windows a decade later
and confound the mixing change with internal variability comparable in size to
the signal. So the experiment is \textbf{seven new runs, not eight}.

\hd{The members.}
Seven perturbations, geometric in powers of two, spanning a factor of $128$
about the reference value:

\smallskip
\begin{center}
\footnotesize
\begin{tabular}{l*{8}{r}}
\toprule
$\kv$ ($\mathrm{m^{2}\,s^{-1}}$)
 & 3.0e$-$6 & 6.0e$-$6 & \textbf{1.2e$-$5} & 2.4e$-$5
 & 4.8e$-$5 & 9.6e$-$5 & 1.9e$-$4 & 3.8e$-$4\\
$\times$ reference
 & 0.25 & 0.5 & \textbf{1 (done)} & 2 & 4 & 8 & 16 & 32\\
\bottomrule
\end{tabular}
\end{center}
\smallskip

Every value lies inside the admissible range already declared for this control.
The spacing is logarithmic because mixing acts on stratification closer to
$\log\kv$ than to \kv, and because the reference sits low in that range, leaving
most of it above.
